%----------------------------------------------------------------------------------------
% FUNCTIES EN HUN EIGENSCHAPPEN. Alle opgaven zijn eigen werk.
%
% Opbouw: Start, leerstof (theorie, werkwijze of uitleg, voorbeeld, korte oefening), en op een
% nieuwe pagina de Opdrachten (door elkaar, stijgende moeilijkheid), Koppel en Test jezelf.
% Labels van dit hoofdstuk beginnen met h1-.
%----------------------------------------------------------------------------------------

\mychapter{1}{Functies en hun eigenschappen}

\lkr{Leerplan B+S': LPD 5-6 (algemene functies). Dit hoofdstuk herhaalt en ordent wat leerlingen uit de 2de graad kennen over eerste- en tweedegraadsfuncties: domein, bereik, nulwaarden, tekenverloop, stijgen en dalen, extrema, en voegt symmetrie toe. Alles wordt eerst op een grafiek gelezen en pas daarna berekend. Afgeleiden komen later (Afgeleiden): extrema worden hier afgelezen, niet berekend.}

Een functie beschrijft hoe een grootheid afhangt van een andere: de temperatuur van het uur, de kostprijs van het aantal stuks, de hoogte van een bal van de tijd. In de 2de graad onderzocht je eerste- en tweedegraadsfuncties. In dit hoofdstuk zetten we de begrippen op een rij die je bij \textit{elke} functie gebruikt: waar een functie bestaat, welke waarden ze aanneemt, waar ze positief is, waar ze stijgt, en of haar grafiek symmetrisch is. Bij elk begrip kijk je eerst op de grafiek, en daarna reken je het na met het voorschrift.

\begin{oefeningen}[Start]
\oef{h1-start-1}

\oef{h1-verken}
\end{oefeningen}

%========================================================================================
\section{Domein en bereik}

\subsection{Functie en grafiek}

Bij een functie stop je een getal $x$ in, en krijg je er een getal $f(x)$ uit. Dat mag nooit dubbelzinnig zijn: bij de temperatuur om 15 uur hoort één temperatuur, niet twee. Wel kan het gebeuren dat er niets uitkomt, zoals bij $\frac10$ of $\sqrt{-4}$.

\begin{theorie}
Een \textbf{\textcolor{accent}{functie}} $f$ van $\R$ naar $\R$ koppelt aan elk reëel getal $x$ \textit{hoogstens één} beeld $f(x)$.

De \textbf{\textcolor{accent}{grafiek}} van $f$ is de verzameling punten $\formule{(x, f(x))}$ waarvoor $f(x)$ bestaat.
\end{theorie}

Omdat elke $x$ hoogstens één beeld heeft, ligt boven elke $x$ hoogstens één punt van de grafiek. Een verticale rechte snijdt de grafiek van een functie dus in hoogstens één punt.

\naast{Zo is de cirkel met middelpunt $(0, 0)$ en straal 2 geen grafiek van een functie. De verticale rechte $x = 1$ snijdt hem in $(1, \sqrt3)$ en in $(1, -\sqrt3)$: bij $x = 1$ zouden twee beelden horen. De parabool $y = x^2 - 2$ is wel een grafiek: elke verticale rechte snijdt ze precies één keer.}{%
\begin{tikzpicture}[x=.42cm,y=.42cm,font=\scriptsize]
  \roosterleeg{-3}{1}{3}{-3}{1}{3}
  \kromme{-2.2}{2.2}{(\x)^2-2}
  \draw[grijs,dashed] (1,-3) -- (1,3);
  \vol{(1,-1)}
\end{tikzpicture}\hspace{2mm}
\begin{tikzpicture}[x=.42cm,y=.42cm,font=\scriptsize]
  \roosterleeg{-3}{1}{3}{-3}{1}{3}
  \draw[accent,thick] (0,0) circle (2);
  \draw[grijs,dashed] (1,-3) -- (1,3);
  \vol{(1,1.732)} \vol{(1,-1.732)}
\end{tikzpicture}}

\subsection{Het domein}

Niet elk getal heeft een beeld. In de 2de graad zag je al dat $\frac1x$ niet bestaat voor $x = 0$, en $\sqrt{x}$ niet voor $x < 0$. Voor je met een functie rekent, wil je weten voor welke getallen ze bestaat.

\begin{theorie}
Het \textbf{\textcolor{accent}{domein}} van $f$ is de verzameling van de getallen die een beeld hebben: \quad $\formule{\dom f = \{x \in \R \mid f(x) \text{ bestaat}\}}$
\end{theorie}

\voorbeeld{}{Bepaal het domein van $f(x) = \dfrac{\sqrt{6 - 2x}}{x + 1}$.\label{vb:h1-domein}}{%
\uitleg{Een vierkantswortel bestaat alleen voor getallen die niet negatief zijn, en een noemer mag niet nul zijn.}
\begin{stappenbreed}
6 - 2x \geqslant 0 \iff x \leqslant 3 & Voorwaarde voor de wortel. \\
x + 1 \neq 0 \iff x \neq -1 & Voorwaarde voor de noemer. \\
\dom f = \left]-\infty, 3\right] \setminus \{-1\} & Beide voorwaarden samen.
\end{stappenbreed}}

\begin{werkwijze}
Zoek in het voorschrift alles wat niet voor elk getal bestaat: een \textbf{noemer} moet verschillend zijn van nul, wat onder een \textbf{vierkantswortel} staat (of een andere even wortel) moet positief of nul zijn. Los elke voorwaarde op en neem de getallen die aan alle voorwaarden voldoen.
\end{werkwijze}

Bepaal het domein altijd met het voorschrift zoals het gegeven is. Wie eerst vereenvoudigt, kan een voorwaarde kwijtspelen.

\voorbeeld{}{Bepaal het domein van $f(x) = \dfrac{x^2 + 3x}{x}$, en teken de grafiek.\label{vb:h1-vereenvoudigen}}{%
\naast{%
\uitleg{De noemer is $x$, dus $0$ valt weg. Pas daarna vereenvoudigen we, om de grafiek te tekenen.}
\begin{stappenvrij}{5.4cm}{4.6cm}
\dom f = \R \setminus \{0\} & De noemer mag niet nul zijn. \\
f(x) = \dfrac{x(x + 3)}{x} = x + 3 & Mag, want $x \neq 0$. \\
f(0) \text{ bestaat niet} & Toch is $0 + 3 = 3$.
\end{stappenvrij}
De grafiek is de rechte $y = x + 3$ \textit{zonder} het punt $(0, 3)$: daar staat een open bolletje. Wie meteen $x + 3$ schrijft, vindt het domein $\R$, en dat is fout.}{%
\begin{tikzpicture}[x=.5cm,y=.5cm,font=\scriptsize]
  \rooster{-3}{1}{2}{-1}{1}{5}
  \draw[accent,thick] (-3,0) -- (2,5);
  \open{(0,3)}
\end{tikzpicture}}}

\nujij{Bepaal het domein van $h(x) = \sqrt{x^2 - 9}$.}{$x^2 - 9 \geqslant 0 \iff x \leqslant -3$ of $x \geqslant 3$, dus $\dom h = \left]-\infty, -3\right] \cup \left[3, +\infty\right[$.}

\begin{oefeningen}
\oef{h1-domein}

\oef{h1-domein-valkuil}
\end{oefeningen}

\subsection{Het bereik}

Het domein zegt welke getallen je in een functie mag stoppen. Het bereik zegt wat eruit kan komen.

\begin{theorie}
Het \textbf{\textcolor{accent}{bereik}} van $f$ is de verzameling van de beelden: \quad $\formule{\ber f = \{f(x) \mid x \in \dom f\}}$
\end{theorie}

Op de grafiek lees je het domein af door de grafiek te projecteren op de $x$-as, en het bereik door ze te projecteren op de $y$-as.

\voorbeeld{}{Lees het domein en het bereik af van de functie $f$ met de grafiek hiernaast.\label{vb:h1-aflezen}}{%
\naast{%
\uitleg{Een vol bolletje hoort bij de grafiek, een open bolletje niet. Projecteer de grafiek op de $x$-as voor het domein en op de $y$-as voor het bereik.}
\begin{stappenvrij}{3.6cm}{3.6cm}
\dom f = \left[-3, 4\right[ & $x = 4$ hoort er niet bij. \\
\ber f = [-2, 3] & Laagste punt $-2$, hoogste $3$.
\end{stappenvrij}
Het open bolletje in $(4, 1)$ verandert niets aan het bereik: de waarde $1$ wordt al bereikt, bijvoorbeeld in $x = -3$.}{%
\begin{tikzpicture}[x=.55cm,y=.55cm,font=\scriptsize]
  \rooster{-3}{1}{4}{-2}{1}{3}
  \draw[accent,thick] (-3,1) -- (-1,3) -- (2,-2) -- (4,1);
  \vol{(-3,1)} \open{(4,1)}
\end{tikzpicture}}}

\begin{oefeningen}
\oef{h1-aflezen}
\end{oefeningen}

Zonder grafiek moet je het bereik berekenen. Dat vraagt twee dingen. Eerst toon je aan dat elk beeld aan een voorwaarde voldoet, bijvoorbeeld $f(x) \geqslant -3$. Daarna toon je aan dat elk getal dat aan die voorwaarde voldoet ook echt een beeld is. Zonder die tweede stap weet je alleen dat het bereik \textit{binnen} $\left[-3, +\infty\right[$ ligt.

\voorbeeld{}{Bepaal het bereik van $f(x) = x^2 + 4x + 1$.\label{vb:h1-bereik}}{%
\naast{%
\uitleg{Splits het kwadraat af: dan zie je de kleinste waarde. Los daarna $f(x) = y$ op om te zien dat elke grotere waarde bereikt wordt.}
\begin{stappenvrij}{4.2cm}{4.8cm}
f(x) = (x + 2)^2 - 3 & Kwadraat afsplitsen. \\
f(x) \geqslant -3 & Want $(x + 2)^2 \geqslant 0$. \\
(x + 2)^2 - 3 = y & Neem een $y \geqslant -3$. \\
x = -2 + \sqrt{y + 3} & Bestaat, want $y + 3 \geqslant 0$. \\
\ber f = \left[-3, +\infty\right[ & Beide stappen samen.
\end{stappenvrij}}{%
\begin{tikzpicture}[x=.5cm,y=.42cm,font=\scriptsize]
  \rooster{-5}{1}{1}{-3}{1}{4}
  \kromme{-4.6}{.6}{(\x+2)^2-3}
  \vol{(-2,-3)}
\end{tikzpicture}}}

\begin{werkwijze}
Zoek een ondergrens of bovengrens voor $f(x)$, bijvoorbeeld door het kwadraat af te splitsen of omdat een wortel nooit negatief is. Toon daarna aan dat elk getal $y$ binnen die grens een beeld is: los $f(x) = y$ op en ga na dat de oplossing in het domein ligt.
\end{werkwijze}

\begin{oefeningen}
\oef{h1-bereik}
\end{oefeningen}

%========================================================================================
\section{Nulwaarden en tekenverloop}

Waar is de winst van een bedrijf positief? Wanneer ligt een bal boven de grond? Zulke vragen gaan over het teken van een functie: is $f(x)$ positief, negatief of nul? Op de grafiek is dat eenvoudig: boven of onder de $x$-as. Met het voorschrift werk je met nulwaarden en een tekentabel.

\begin{theorie}
Een \textbf{\textcolor{accent}{nulwaarde}} van $f$ is een getal $a \in \dom f$ met $\formule{f(a) = 0}$. Het punt $(a, 0)$ is een snijpunt van de grafiek met de $x$-as.

Het \textbf{\textcolor{accent}{tekenverloop}} van $f$ zegt waar $f(x) > 0$ (grafiek boven de $x$-as) en waar $f(x) < 0$ (grafiek eronder). Je noteert het in een \textbf{\textcolor{accent}{tekentabel}}.
\end{theorie}

\subsection{Het teken tussen twee nulwaarden}

Bij een tekentabel bereken je meestal één functiewaarde per interval tussen twee nulwaarden. Waarom volstaat dat? Omdat het teken daar niet kan veranderen, als je de functie ontbonden hebt in factoren waarvan je het teken kent.

\begin{theorie}
\begin{stelling}{Teken tussen twee nulwaarden}
Is $f(x)$ een product van een constante, eerstegraadsfactoren en tweedegraadsfactoren met discriminant $D < 0$, dan heeft $f$ hetzelfde teken op elk interval dat geen nulwaarde van $f$ bevat.
\end{stelling}

\begin{bewijs}
Neem een interval $I$ zonder nulwaarde van $f$. Een product is nul als een factor nul is, dus ook geen enkele factor is nul op $I$.

Een constante factor heeft vanzelf één teken. Een eerstegraadsfactor $mx + q$ heeft één nulwaarde $-\frac{q}{m}$, en heeft links ervan het ene teken en rechts ervan het andere. Die nulwaarde ligt niet in $I$, dus $I$ ligt er helemaal links of helemaal rechts van, en de factor heeft één teken op $I$.

Een tweedegraadsfactor $ax^2 + bx + c$ met $D < 0$ heeft geen nulwaarden en heeft overal het teken van $a$, want
\bewijsregel{ax^2 + bx + c = a\Bigl(\bigl(x + \tfrac{b}{2a}\bigr)^2 - \tfrac{D}{4a^2}\Bigr) \quad\text{met}\quad \bigl(x + \tfrac{b}{2a}\bigr)^2 - \tfrac{D}{4a^2} > 0 .}
Elke factor heeft dus één teken op $I$, en volgens de tekenregel heeft hun product dat ook.
\end{bewijs}
\end{theorie}

De stelling zegt: een functie die je zo kunt ontbinden, verandert alleen van teken in een nulwaarde. Daarom volstaat één proefgetal per interval. Elke veeltermfunctie kan je in principe zo ontbinden, maar dat bewijs valt buiten deze cursus. In de praktijk werk je met veeltermen die je zelf kunt ontbinden.

\voorbeeld{}{Bepaal de nulwaarden en het tekenverloop van $f(x) = x^3 - 4x$.\label{vb:h1-teken}}{%
\uitleg{Ontbind in factoren. Tussen twee opeenvolgende nulwaarden verandert het teken niet, dus één getal in elk interval volstaat.}
\begin{stappenbreed}
f(x) = x(x^2 - 4) = x(x - 2)(x + 2) & Gemeenschappelijke factor, dan merkwaardig product. \\
\text{nulwaarden: } -2,\ 0,\ 2 & Een product is nul als een factor nul is. \\
f(-3) = -15,\ f(-1) = 3,\ f(1) = -3,\ f(3) = 15 & Eén getal in elk interval.
\end{stappenbreed}
\figuren{\tekentabel{cccccccc}{x & & -2 & & 0 & & 2 & }{f(x) & - & 0 & + & 0 & - & 0 & +}}
Dus $f(x) > 0$ voor $x \in \left]-2, 0\right[ \cup \left]2, +\infty\right[$, en $f(x) < 0$ voor $x \in \left]-\infty, -2\right[ \cup \left]0, 2\right[$.}

\begin{oefeningen}
\oef{h1-teken-factoren}
\end{oefeningen}

\subsection{Een factor met een exponent}

Staat een factor meer dan één keer in het product, dan wisselt het teken in die nulwaarde niet altijd. Bij $(x - 1)^2$ zie je dat meteen: een kwadraat is nooit negatief, links noch rechts van $1$.

\begin{theorie}
\begin{stelling}{Even of oneven exponent}
Is $f(x) = (x - a)^n \cdot g(x)$, met $g(x)$ een product zoals in de vorige stelling en $g(a) \neq 0$, dan wisselt $f$ van teken in $a$ als $n$ oneven is, en niet als $n$ even is.
\end{stelling}

\begin{bewijs}
$g$ heeft maar eindig veel nulwaarden, en $a$ is er geen van. Neem dus een interval rond $a$ zonder nulwaarde van $g$: daarop heeft $g(x)$ één teken, volgens de vorige stelling. Is $n$ even, dan is $(x - a)^n > 0$ voor elke $x \neq a$, links en rechts van $a$: het teken van $f$ is dat van $g$, aan beide kanten hetzelfde. Is $n$ oneven, dan heeft $(x - a)^n$ het teken van $x - a$: negatief links van $a$, positief rechts ervan. Het teken van $f$ is dan links het tegengestelde van rechts.
\end{bewijs}
\end{theorie}

Op de grafiek: bij een oneven exponent snijdt de grafiek de $x$-as, bij een even exponent raakt ze de $x$-as en keert ze terug.

\voorbeeld{}{Stel de tekentabel op van $f(x) = x^3(x - 2)^2(x + 1)$.\label{vb:h1-exponent}}{%
\naast{%
\uitleg{De nulwaarden en hun exponenten staan al in het voorschrift. Bepaal het teken rechts van de grootste nulwaarde, en ga dan naar links.}
\begin{stappenvrij}{3.9cm}{3.8cm}
-1,\ 0,\ 2 & De nulwaarden. \\
f(3) = 27 \cdot 1 \cdot 4 > 0 & Rechts van $2$: positief. \\
(x - 2)^2 & Even: in $2$ geen wissel. \\
x^3 & Oneven: in $0$ wel. \\
x + 1 & Oneven: in $-1$ wel.
\end{stappenvrij}}{%
\begin{tikzpicture}[x=.75cm,y=.5cm,font=\scriptsize]
  \rooster{-1}{1}{2}{-1}{1}{4}
  \kromme{-1.2}{2.3}{(\x)^3*(\x-2)^2*(\x+1)}
\end{tikzpicture}}
\figuren{\tekentabel{cccccccc}{x & & -1 & & 0 & & 2 & }{f(x) & + & 0 & - & 0 & + & 0 & +}}
Op de grafiek snijdt $f$ de $x$-as in $-1$ en $0$, en raakt ze de $x$-as in $2$.}

\begin{oefeningen}
\oef{h1-teken-veelterm}
\end{oefeningen}

\subsection{Een functie met een noemer}

Bestaat $f$ niet in een getal, dan kan het teken daar ook wisselen. Een quotiënt heeft volgens de tekenregel hetzelfde teken als het product van teller en noemer. Beide stellingen gelden dus ook voor een breuk, als je de nulwaarden van de noemer erbij neemt. Zet zo'n getal mee in de tabel, met een streep: daar bestaat $f$ niet.

\voorbeeld{}{Stel de tekentabel op van $f(x) = \dfrac{x - 1}{x + 2}$.\label{vb:h1-teken-breuk}}{%
\uitleg{Teller en noemer krijgen elk een eigen rij; het teken van het quotiënt volgt uit de tekenregel.}
\figuren{\tekentabel{cccccc}{x & & -2 & & 1 & \\\hline x - 1 & - & - & - & 0 & + \\ x + 2 & - & 0 & + & + & +}{f(x) & + & | & - & 0 & +}}
Nulwaarde $1$; in $-2$ bestaat $f$ niet.}

\begin{werkwijze}
Ontbind het voorschrift in factoren en bepaal de nulwaarden van elke factor (ook van de noemer). Zet die getallen in volgorde bovenaan de tekentabel. Bepaal het teken in elk interval: met een rij per factor en de tekenregel, of door een getal in te vullen en de exponenten te volgen.
\end{werkwijze}

\begin{oefeningen}
\oef{h1-teken}
\end{oefeningen}

%========================================================================================
\section{Stijgen, dalen en extrema}

Een grafiek die je van links naar rechts volgt, gaat omhoog of omlaag. Dat heet stijgen en dalen. Waar ze van richting verandert, ligt een top of een dal: een extremum. In deze sectie maken we die woorden nauwkeurig, zodat je ze ook kunt aantonen.

\subsection{Stijgen en dalen}

\naast{Een functie stijgt als een grotere $x$ altijd een groter beeld geeft. Dat moet gelden voor \textit{elke} twee getallen van het interval, niet alleen voor een paar die je toevallig kiest. In de figuur ligt $x_2$ rechts van $x_1$, en daarom ligt ook het punt van de grafiek boven $x_2$ hoger.}{%
\begin{tikzpicture}[x=.7cm,y=.5cm,font=\scriptsize]
  \draw[->,grijs] (-.2,0) -- (4.3,0) node[right,text=black] {$x$};
  \draw[->,grijs] (0,-.2) -- (0,4.3) node[above,text=black] {$y$};
  \kromme{.4}{4}{.6+.25*(\x)^2-.02*(\x)^3}
  \draw[grijs,dashed] (1.2,0) node[below,text=black] {$x_1$} |- (0,{.6+.25*1.44-.02*1.728}) node[left,text=black] {$f(x_1)$};
  \draw[grijs,dashed] (3.2,0) node[below,text=black] {$x_2$} |- (0,{.6+.25*10.24-.02*32.768}) node[left,text=black] {$f(x_2)$};
  \vol{(1.2,{.6+.25*1.44-.02*1.728})} \vol{(3.2,{.6+.25*10.24-.02*32.768})}
\end{tikzpicture}}

\begin{theorie}
Een functie $f$ is \textbf{\textcolor{accent}{stijgend}} op een interval $I$ als voor alle $x_1, x_2 \in I$ geldt: $\formule{x_1 < x_2 \Rightarrow f(x_1) < f(x_2)}$, en \textbf{\textcolor{accent}{dalend}} als $\formule{x_1 < x_2 \Rightarrow f(x_1) > f(x_2)}$.
\end{theorie}

Uit de 2de graad weet je dat $f(x) = mx + q$ stijgt als $m > 0$ en daalt als $m < 0$. Voor andere functies kun je het met de definitie aantonen: neem twee willekeurige getallen $x_1 < x_2$ en bekijk het verschil $f(x_2) - f(x_1)$.

\voorbeeld{}{Toon met de definitie aan dat $f(x) = \sqrt{x}$ stijgend is op $\left[0, +\infty\right[$.\label{vb:h1-stijgend}}{%
\uitleg{Een verschil van twee wortels is moeilijk te schatten. Vermenigvuldig teller en noemer met de som van de wortels: dan verdwijnen de wortels uit de teller.}
\begin{stappenbreed}
0 \leqslant x_1 < x_2 & Twee willekeurige getallen uit het interval. \\
\sqrt{x_2} - \sqrt{x_1} = \dfrac{(\sqrt{x_2} - \sqrt{x_1})(\sqrt{x_2} + \sqrt{x_1})}{\sqrt{x_2} + \sqrt{x_1}} = \dfrac{x_2 - x_1}{\sqrt{x_2} + \sqrt{x_1}} & De noemer is niet nul, want $x_2 > x_1 \geqslant 0$. Merkwaardig product in de teller. \\
\dfrac{x_2 - x_1}{\sqrt{x_2} + \sqrt{x_1}} > 0 & Teller en noemer zijn positief. \\
\sqrt{x_1} < \sqrt{x_2} & Dus $f$ is stijgend.
\end{stappenbreed}}

\begin{werkwijze}
Neem twee willekeurige getallen $x_1 < x_2$ uit het interval. Schrijf het verschil $f(x_2) - f(x_1)$ zo dat je het teken ziet, bijvoorbeeld door te ontbinden. Is het altijd positief, dan stijgt $f$; is het altijd negatief, dan daalt $f$.
\end{werkwijze}

\nujij{Toon met de definitie aan dat $f(x) = 5 - 2x$ dalend is op $\R$.}{Neem $x_1 < x_2$. Dan is $f(x_2) - f(x_1) = -2(x_2 - x_1) < 0$, dus $f(x_1) > f(x_2)$.}

\begin{oefeningen}
\oef{h1-stijgend-def}
\end{oefeningen}

\subsection{Extrema}

Waar een grafiek van stijgen naar dalen gaat, ligt een top. Die hoeft niet het hoogste punt van de hele grafiek te zijn: het volstaat dat er in de buurt niets hoger ligt.

\begin{theorie}
$f$ bereikt een \textbf{\textcolor{accent}{maximum}} $f(a)$ in $a$ als $\formule{f(x) \leqslant f(a)}$ voor alle $x$ uit het domein in de buurt van $a$, en een \textbf{\textcolor{accent}{minimum}} als $\formule{f(x) \geqslant f(a)}$. Maxima en minima heten samen \textbf{\textcolor{accent}{extrema}}.

\begin{stelling}{Van stijgen naar dalen}
Is $f$ stijgend op $[b, a]$ en dalend op $[a, c]$, dan is $f(a)$ een maximum: $\formule{f(x) \leqslant f(a)}$ voor alle $x \in [b, c]$.

Is $f$ dalend op $[b, a]$ en stijgend op $[a, c]$, dan is $f(a)$ een minimum.
\end{stelling}

\begin{bewijs}
Neem $x \in [b, c]$. Is $x < a$, dan liggen $x$ en $a$ in $[b, a]$, waar $f$ stijgt, dus $f(x) < f(a)$. Is $x > a$, dan liggen ze in $[a, c]$, waar $f$ daalt, dus $f(a) > f(x)$. Is $x = a$, dan is $f(x) = f(a)$. In elk geval is $f(x) \leqslant f(a)$. Voor een minimum keren alle ongelijkheden om.
\end{bewijs}
\end{theorie}

Je vat het verloop samen in een \textbf{verlooptabel}: een pijl omhoog waar $f$ stijgt, omlaag waar $f$ daalt, met de extrema erboven of eronder. Volgens de stelling ligt tussen een pijl omhoog en een pijl omlaag een maximum, en tussen een pijl omlaag en een pijl omhoog een minimum.

\voorbeeld{}{De grafiek hiernaast is die van $f(x) = x^3 - 3x$. Lees af waar $f$ stijgt en daalt, en geef de extrema.\label{vb:h1-verloop}}{%
\naast{%
\uitleg{Volg de grafiek van links naar rechts. Waar ze van stijgen naar dalen gaat, ligt een maximum; omgekeerd een minimum.}
\begin{stappenvrij}{4.2cm}{3.4cm}
\text{stijgend op } \left]-\infty, -1\right] & Tot de top. \\
\text{dalend op } [-1, 1] & Van top naar dal. \\
\text{stijgend op } \left[1, +\infty\right[ & Daarna weer omhoog. \\
\text{max. } f(-1) = 2 & De top. \\
\text{min. } f(1) = -2 & Het dal.
\end{stappenvrij}}{%
\begin{tikzpicture}[x=.6cm,y=.42cm,font=\scriptsize]
  \rooster{-2}{1}{2}{-4}{1}{4}
  \kromme{-2.2}{2.2}{(\x)^3-3*\x}
  \vol{(-1,2)} \vol{(1,-2)}
\end{tikzpicture}}
\figuren{\tekentabel{cccccc}{x & & -1 & & 1 & }{f(x) & \nearrow & 2 & \searrow & -2 & \nearrow}}
Het maximum $2$ is niet het hoogste punt van de grafiek: voor grote $x$ wordt $f(x)$ groter dan $2$, en voor kleine $x$ kleiner dan $-2$.}

\begin{werkwijze}
Lees de grafiek van links naar rechts en noteer de $x$-waarden waar ze van richting verandert. Geef de intervallen waar ze stijgt en daalt, en per top of dal het extremum als \textit{functiewaarde} $f(a)$ en de plaats $a$ waar het bereikt wordt.
\end{werkwijze}

\nujij{Geef het verloop en de extrema van de functie uit voorbeeld~\ref{vb:h1-aflezen}.}{Stijgend op $[-3, -1]$, dalend op $[-1, 2]$, stijgend op $\left[2, 4\right[$. Maximum $3$ in $x = -1$ (absoluut), minimum $-2$ in $x = 2$ (absoluut).}

\begin{oefeningen}
\oef{h1-verloop}

\needspace{8\baselineskip}
\oef{h1-verlooptabel}
\end{oefeningen}

\subsection{Absolute extrema}

Soms wil je de grootste waarde van een functie op heel haar domein: de grootste winst, de hoogste temperatuur van de dag. Dat is een strengere eis dan een top in de buurt.

\begin{theorie}
Een maximum $f(a)$ is \textbf{\textcolor{accent}{absoluut}} als $\formule{f(x) \leqslant f(a)}$ voor heel het domein. Een minimum $f(a)$ is absoluut als $\formule{f(x) \geqslant f(a)}$ voor heel het domein.
\end{theorie}

Een extremum kan ook op de rand van het domein liggen. Daar kijk je maar naar één kant.

\voorbeeld{}{Geef de extrema van $f(x) = x^2$ met domein $[-1, 2]$. Welke zijn absoluut?\label{vb:h1-absoluut}}{%
\naast{%
\uitleg{Bekijk de grafiek: $f$ daalt op $[-1, 0]$ en stijgt op $[0, 2]$. Kijk ook naar de twee randpunten.}
\begin{stappenvrij}{3.4cm}{6.2cm}
\text{min. } f(0) = 0 & Van dalen naar stijgen; absoluut, want $x^2 \geqslant 0$. \\
\text{max. } f(-1) = 1 & Rechts ervan daalt $f$; niet absoluut. \\
\text{max. } f(2) = 4 & Links ervan stijgt $f$; absoluut.
\end{stappenvrij}}{%
\begin{tikzpicture}[x=.6cm,y=.42cm,font=\scriptsize]
  \rooster{-1}{1}{2}{0}{1}{4}
  \kromme{-1}{2}{(\x)^2}
  \vol{(-1,1)} \vol{(0,0)} \vol{(2,4)}
\end{tikzpicture}}
Met domein $\left]-1, 2\right[$ heeft $f$ geen enkel maximum meer: vlak naast elke $x$ van het domein ligt een getal met een groter beeld. De beelden komen willekeurig dicht bij $4$, maar $4$ wordt niet bereikt.}

\begin{oefeningen}
\oef{h1-absoluut}
\end{oefeningen}

%========================================================================================
\section{Even en oneven functies}

Veel grafieken zijn symmetrisch, zoals die van $x^2$ en $x^3$. Dat is handig: ken je de helft van de grafiek, dan ken je ook de andere helft. In deze sectie leer je die symmetrie herkennen aan het voorschrift.

\subsection{Even functies}

Bij $f(x) = x^2$ hebben $3$ en $-3$ hetzelfde beeld, $9$. Een functie waarbij dat voor elk getal en zijn tegengestelde geldt, heet even.

\begin{theorie}
Een functie $f$ is \textbf{\textcolor{accent}{even}} als voor elke $x \in \dom f$ ook $-x \in \dom f$ en $\formule{f(-x) = f(x)}$.

\begin{stelling}{Symmetrie om de \mbox{$y$-as}}
Een functie is even als en slechts als haar grafiek symmetrisch is ten opzichte van de \mbox{$y$-as}.
\end{stelling}

\begin{bewijs}
Het spiegelbeeld van het punt $(a, b)$ om de $y$-as is $(-a, b)$.

Stel dat $f$ even is, en dat $(a, b)$ op de grafiek ligt: $a \in \dom f$ en $b = f(a)$. Dan is $-a \in \dom f$ en $f(-a) = f(a) = b$, dus ook $(-a, b)$ ligt op de grafiek.

Stel omgekeerd dat met elk punt $(a, f(a))$ ook $(-a, f(a))$ op de grafiek ligt. Dan bestaat $f(-a)$, dus $-a \in \dom f$, en $f(-a) = f(a)$.
\end{bewijs}
\end{theorie}

Het domein van een even functie is dus zelf symmetrisch: met $3$ hoort ook $-3$ erbij. Vergeet je dat, dan lijkt een functie soms even terwijl ze het niet is.

\voorbeeld{}{Is $f(x) = x^2$ met domein $[-1, 3]$ een even functie?\label{vb:h1-domein-symm}}{%
\naast{%
\uitleg{De gelijkheid $(-x)^2 = x^2$ klopt, maar kijk eerst naar het domein.}
\begin{stappenvrij}{3.4cm}{4.2cm}
2 \in \dom f & $2$ ligt in $[-1, 3]$. \\
-2 \notin \dom f & $-2$ ligt er niet in. \\
f \text{ is niet even} & De eerste voorwaarde faalt.
\end{stappenvrij}
De grafiek is geen symmetrische parabool: rechts loopt ze verder dan links.}{%
\begin{tikzpicture}[x=.5cm,y=.36cm,font=\scriptsize]
  \rooster{-3}{1}{3}{0}{1}{9}
  \kromme{-1}{3}{(\x)^2}
  \vol{(-1,1)} \vol{(3,9)}
\end{tikzpicture}}}

\subsection{Oneven functies}

Bij $f(x) = x^3$ hebben $2$ en $-2$ tegengestelde beelden: $8$ en $-8$. Een functie waarbij dat voor elk getal geldt, heet oneven.

\begin{theorie}
Een functie $f$ is \textbf{\textcolor{accent}{oneven}} als voor elke $x \in \dom f$ ook $-x \in \dom f$ en $\formule{f(-x) = -f(x)}$.

\begin{stelling}{Symmetrie om de oorsprong}
Een functie is oneven als en slechts als haar grafiek puntsymmetrisch is ten opzichte van de oorsprong.
\end{stelling}

\begin{bewijs}
Het spiegelbeeld van het punt $(a, b)$ om de oorsprong is $(-a, -b)$.

Stel dat $f$ oneven is, en dat $(a, b)$ op de grafiek ligt. Dan is $-a \in \dom f$ en $f(-a) = -f(a) = -b$, dus ook $(-a, -b)$ ligt op de grafiek.

Stel omgekeerd dat met elk punt $(a, f(a))$ ook $(-a, -f(a))$ op de grafiek ligt. Dan is $-a \in \dom f$ en $f(-a) = -f(a)$.
\end{bewijs}
\end{theorie}

\figuren{%
\begin{tikzpicture}[x=.55cm,y=.4cm,font=\scriptsize]
  \roosterleeg{-3}{1}{3}{-4}{1}{3}
  \kromme{-2.9}{2.9}{(\x)^4/4-2*(\x)^2}
  \node[below,text=black] at (0,-4.3) {even: $\frac14x^4 - 2x^2$};
\end{tikzpicture}\hspace{1.2cm}
\begin{tikzpicture}[x=.55cm,y=.4cm,font=\scriptsize]
  \roosterleeg{-3}{1}{3}{-4}{1}{3}
  \kromme{-2.2}{2.2}{(\x)^3-3*\x}
  \vol{(-1,2)} \vol{(1,-2)}
  \draw[accent,dashed] (-1,2) -- (1,-2);
  \node[below,text=black] at (0,-4.3) {oneven: $x^3 - 3x$};
\end{tikzpicture}}

Bij een even functie zie je de symmetrie als een spiegel langs de $y$-as. Bij een oneven functie draai je de grafiek een halve slag om de oorsprong: de top in $(-1, 2)$ komt op het dal in $(1, -2)$.

\voorbeeld{}{Toon aan dat $f(x) = x^3 - 3x$ oneven is.\label{vb:h1-even}}{%
\uitleg{Het domein is $\R$, dus met $x$ zit ook $-x$ erin. Bereken $f(-x)$ en vergelijk met $-f(x)$.}
\begin{stappenbreed}
f(-x) = (-x)^3 - 3(-x) = -x^3 + 3x & Vul $-x$ in. \\
-f(x) = -(x^3 - 3x) = -x^3 + 3x & Hetzelfde: $f(-x) = -f(x)$.
\end{stappenbreed}}

Een oneven functie gaat altijd door de oorsprong, als $0$ in haar domein ligt. Daarmee kun je snel uitsluiten dat een functie oneven is.

\begin{theorie}
\begin{stelling}{Oneven functie in $0$}
Is $f$ oneven en $0 \in \dom f$, dan is $\formule{f(0) = 0}$.
\end{stelling}

\begin{bewijs}
Neem $x = 0$ in de definitie: $f(0) = f(-0) = -f(0)$, dus $2f(0) = 0$ en $f(0) = 0$.
\end{bewijs}
\end{theorie}

\voorbeeld{}{De grafiek van $f$ met domein $[-3, 3]$ is gegeven voor $x \geqslant 0$: de gebroken lijn door $(0, 0)$, $(2, 2)$ en $(3, 1)$. Teken de volledige grafiek als $f$ oneven is.\label{vb:h1-spiegel}}{%
\naast{%
\uitleg{Spiegel elk hoekpunt om de oorsprong: $(a, b) \mapsto (-a, -b)$. Controleer eerst of $f(0) = 0$, anders kan $f$ niet oneven zijn.}
\begin{stappenvrij}{3.4cm}{4.2cm}
f(0) = 0 & Dat klopt. \\
(2, 2) \mapsto (-2, -2) & Hoekpunt spiegelen. \\
(3, 1) \mapsto (-3, -1) & Eindpunt spiegelen.
\end{stappenvrij}
Het gegeven deel is grijs, het gespiegelde deel in kleur.}{%
\begin{tikzpicture}[x=.45cm,y=.45cm,font=\scriptsize]
  \rooster{-3}{1}{3}{-2}{1}{2}
  \draw[grijs,thick] (0,0) -- (2,2) -- (3,1);
  \draw[accent,thick] (0,0) -- (-2,-2) -- (-3,-1);
  \vol{(-3,-1)} \fill[grijs] (3,1) circle (2.2pt);
\end{tikzpicture}}}

\begin{oefeningen}
\oef{h1-symm-domein}

\needspace{8\baselineskip}
\oef{h1-symmetrie}
\end{oefeningen}

\subsection{Even en oneven machten}

Waarom heten deze functies even en oneven? Omdat machten met een even exponent even functies zijn, en machten met een oneven exponent oneven: $(-x)^2 = x^2$, maar $(-x)^3 = -x^3$. Dat geldt ook voor sommen van zulke machten.

\begin{theorie}
\begin{stelling}{Even en oneven machten}
Een veeltermfunctie met alleen even exponenten is even. Een veeltermfunctie met alleen oneven exponenten is oneven.
\end{stelling}

\begin{bewijs}
Voor een natuurlijk getal $n$ is $(-x)^n = (-1)^n x^n$, en dat is $x^n$ als $n$ even is, en $-x^n$ als $n$ oneven is. Het domein is $\R$. Heeft $f(x) = a_k x^k + \ldots + a_1 x + a_0$ alleen even exponenten, dan verandert geen enkele term als je $x$ door $-x$ vervangt, dus $f(-x) = f(x)$. Heeft $f$ alleen oneven exponenten, dan verandert elke term van teken, dus $f(-x) = -f(x)$.
\end{bewijs}
\end{theorie}

Let op: een constante term, zoals $+3$, is $3x^0$, en $0$ is een even exponent. Zo is $x^4 - x^2 + 3$ even, maar $x^3 + 1$ niet oneven: daar staan een oneven en een even exponent door elkaar. Staan beide soorten door elkaar, dan toon je met één tegenvoorbeeld aan dat de functie even noch oneven is.

\voorbeeld{}{Toon aan dat $g(x) = x^2 + x$ even noch oneven is.\label{vb:h1-geen}}{%
\uitleg{Er staat een even en een oneven exponent: de stelling zegt dus niets. Zoek één getal waarvoor de gelijkheid niet klopt.}
\begin{stappenbreed}
g(1) = 2,\ g(-1) = 0 & Twee beelden berekenen. \\
g(-1) \neq g(1) \text{ en } g(-1) \neq -g(1) & Dus niet even en niet oneven.
\end{stappenbreed}}

Eén tegenvoorbeeld volstaat om een algemene uitspraak te weerleggen. Omgekeerd volstaan een paar voorbeelden nooit om aan te tonen dat een functie wel even is: daarvoor reken je $f(-x)$ uit met de letter $x$.

\begin{oefeningen}
\oef{h1-even}
\end{oefeningen}

%========================================================================================
\opdrachtenkop

\begin{oefeningen}
\needspace{12\baselineskip}
\oef{h1-o1}

\oef{h1-o2}

\oef{h1-o3}

\oef{h1-o12}

\needspace{12\baselineskip}
\oef{h1-o10}

\needspace{8\baselineskip}
\oef{h1-o11}

\needspace{10\baselineskip}
\oef{h1-o14}

\needspace{9\baselineskip}
\oef{h1-o16}

\needspace{8\baselineskip}
\oef{h1-o17}

\needspace{8\baselineskip}
\oef{h1-o4}

\needspace{9\baselineskip}
\oef{h1-o13}

\oef{h1-o5}

\needspace{10\baselineskip}
\oef{h1-o15}

\needspace{8\baselineskip}
\oef{h1-o20}

\needspace{9\baselineskip}
\oef{h1-o18}

\needspace{9\baselineskip}
\oef{h1-o26}

\oef{h1-o6}

\needspace{9\baselineskip}
\oef{h1-o21}

\oef{h1-o7}

\needspace{10\baselineskip}
\oef{h1-o19}

\needspace{8\baselineskip}
\oef{h1-o22}

\oef{h1-o9}

\needspace{9\baselineskip}
\oef{h1-o8}

\needspace{8\baselineskip}
\oef{h1-o24}

\needspace{9\baselineskip}
\oef{h1-o25}

\oef{h1-o23}
\end{oefeningen}

%----------------------------------------------------------------------------------------
% KOPPEL: verbanden met andere leerstof
%----------------------------------------------------------------------------------------
\needspace{20\baselineskip}
\opdrachtsectie{Koppel}

\lkr{Statistiek en Integralen: kijk na wat al gezien is. Het pakket legt de volgorde niet vast.}

\begin{oefeningen}
\oef{h1-koppel-1}

\oef{h1-koppel-3}
\end{oefeningen}

%----------------------------------------------------------------------------------------
\opdrachtsectie{Test jezelf}

\begin{oefeningen}
\oef{h1-t1}

\oef{h1-t5}

\oef{h1-t2}

\oef{h1-t6}

\oef{h1-t3}

\oef{h1-t4}

\oef{h1-t7}
\end{oefeningen}

\herhaalschema{%
  \fv{oef:h1-t1}{vb:h1-domein}{oef:h1-domein}%
  \fv{oef:h1-t5}{vb:h1-aflezen}{oef:h1-aflezen}%
  \fv{oef:h1-t2}{vb:h1-teken}{oef:h1-teken-factoren}%
  \fv{oef:h1-t6}{vb:h1-teken-breuk}{oef:h1-teken}%
  \fv{oef:h1-t3}{vb:h1-verloop}{oef:h1-verloop}%
  \fv{oef:h1-t4}{vb:h1-even}{oef:h1-even}%
  \fv{oef:h1-t7}{vb:h1-spiegel}{oef:h1-symmetrie}}

\ifzelfstudie
\section*{Maak je eigen samenvatting}
Sluit dit boek en schrijf in je schrift: wat een functie, het domein, het bereik en een nulwaarde zijn, hoe je een tekentabel opstelt, de definities van stijgend, dalend, maximum en minimum, en wanneer een functie even of oneven is (met de symmetrie van de grafiek). Vergelijk daarna met de kaders.
\fi
